%% Section body, shared verbatim by ../paper/main.tex (preprint) and
%% ../submission/body.tex (Information Sciences submission).  The section heading and
%% label live in the including file, because the two versions number sections the same
%% way but the submission adds the journal's required unnumbered sections after them.
Many estimation problems now arrive as a stream: observations appear one at a time or in
small batches, they are too numerous to store, and a decision or a report is required at
any moment. Stochastic approximation \citep{robbins1951stochastic} is the natural tool,
and averaged stochastic gradient descent \citep{ruppert1988efficient,polyak1992acceleration}
is asymptotically efficient with $O(d)$ work per observation. Its weakness is
conditioning: when the curvature of the objective has eigenvalues spread over orders of
magnitude, first-order methods crawl \citep{bottou2018optimization,leluc2023asymptotic}.
Second-order methods fix the conditioning but appear to cost $O(d^3)$ per step.

A line of work on \emph{inversion-free} stochastic Newton methods removes that obstacle.
When the Hessian has the form
$\nabla^2F(\theta)=\E[\alpha(\theta;\xi)\,\Psi(\theta;\xi)\Psi(\theta;\xi)^{\!\top}]$---as
it does for linear, logistic, softmax, Poisson and ridge regression, and for the
geometric median---the running curvature estimate is a sum of rank-one terms, so its
inverse can be maintained directly by the Sherman--Morrison rank-one update
\citep{sherman1950adjustment} --- also called the Riccati form, the name we use in the proofs
--- without ever forming or inverting a matrix
\citep{bercu2020efficient,cenac2025efficient,boyer2023asymptotic,godichonbaggioni2024online}.
\citet{godichonbaggioni2025adaptive} push this to its natural conclusion: thinning the
rank-one updates and taking mini-batches of size $d$ gives an estimator whose \emph{total}
cost is $O(dN)$---first-order cost---while remaining asymptotically efficient. A
concurrent construction achieves the same budget by masking Hessian columns
\citep{godichonbaggioni2026complexity}.

All of this theory assumes the stream is independent and identically distributed. The
streams that motivate it are not. Hourly traffic volume, pollutant concentrations,
queue lengths, and user activity are serially dependent, and after any regression
adjustment the \emph{scores} remain serially dependent: in the two real streams we study,
the lag-one residual autocorrelation after a fixed weather-and-calendar adjustment is
$0.76$ (traffic) and $0.92$ (air quality). The consequence is not a small correction. Writing $g_i=\nabla_\theta f(\ts;\xi_i)$ for the score at the truth,
$\Gamma_k=\Cov(g_0,g_k)$, $H=\nabla^2F(\ts)$ and
\begin{equation}
S \;=\; \sum_{k\in\mathbb{Z}}\Gamma_k
\qquad\text{(the \emph{long-run} score covariance),}
\label{eq:S}
\end{equation}
any efficient streaming $M$-estimator satisfies
$\sqrt{N}(\hat\theta_N-\ts)\rightsquigarrow \mathcal{N}(0,H^{-1}SH^{-1})$, whereas the
independent-data theory certifies $\mathcal{N}(0,H^{-1}\Gamma_0H^{-1})$. The ratio of the
two variances, coordinatewise, is a factor $\kappa_j$ that does not shrink with $N$. It is
not a rounding error: we estimate it at about $2$ for the traffic stream and about $20$ for
the air-quality stream, and at $\kappa_j=3$ a nominal $95\%$ interval covers $75\%$ of the
time \emph{forever}, however long the stream runs.

\paragraph{This paper.}
We ask what has to change in a streaming inversion-free quasi-Newton method so that both
its point estimate and its \emph{uncertainty} are correct on a dependent stream, without
leaving the $O(dN)$ budget. Our answer is that one structural choice does both jobs.
Consume the stream in \emph{contiguous blocks} of length $b$ and take one preconditioned
step per block; place curvature updates on a \emph{gap} of $m$ observations. Then

\begin{itemize}[leftmargin=1.4em,itemsep=1pt,topsep=2pt]
\item the $b$-averaged gradients the algorithm forms anyway are, up to a factor $b$, the
      non-overlapping batch means of the score process, so they \emph{are} an estimator of
      $S$---the statistically correct object and the computational shortcut coincide; and
\item gapping decorrelates the rank-one curvature terms, which is what makes the
      inversion-free recursion behave, at matched cost, as it would on independent data.
\end{itemize}

\noindent
Neither device substitutes for the other: gapping alone leaves the variance wrong, and
blocking alone leaves the curvature recursion aggregating dependent terms. We call the
resulting estimator \bgsn{} (blocked--gapped streaming Newton).

\paragraph{Contributions.}
\begin{enumerate}[leftmargin=1.6em,itemsep=2pt,topsep=2pt]
\item \textbf{Theory for the inversion-free curvature recursion under dependence}
      (\Cref{sec:theory}). We prove almost sure convergence, the rate
      $\norm{\theta_T-\ts}^2=O(\log N/N)$, and a central limit theorem with the long-run
      sandwich $H^{-1}SH^{-1}$, for a rank-one Sherman--Morrison curvature recursion with
      gapped updates on a geometrically mixing stream
      (\Cref{thm:curv,thm:as,thm:rate,thm:clt}). The independent-data analysis rests on
      the score being a martingale difference sequence with respect to the algorithm's
      filtration; dependence destroys that, and existing dependent-data results are
      first-order, while existing streaming second-order inference assumes conditionally
      unbiased gradients. A by-product (\Cref{prop:blockfree}) is that for this algorithm
      blocking is \emph{free}: the leading term of the error is
      $-N^{-1}H^{-1}\sum_{i\le N}g_i$ for \emph{every} $b$, so the asymptotic variance does
      not depend on the block length.
\item \textbf{A long-run covariance estimator from block gradients, and why it is fast}
      (\Cref{sec:lrv}). $S$ is estimated at $O(d^2)$ per block and $O(d^2)$ memory from
      the same block gradients. The key point is \emph{what} is being blocked: batch means
      taken over the \emph{iterate} path need blocks long enough for the iterates to mix,
      which forces $b\asymp N^{3/4}$ and caps the rate at $N^{-1/8}$
      \citep{zhu2023online,roy2023online,samsonov2025statistical}; block \emph{gradients}
      inherit their dependence from the data, so $b\asymp\log N$ suffices and
      \Cref{thm:lrv} gives $\widetilde O(N^{-1/2})$. This does not contradict the
      minimax rate $\Theta(N^{-(1-a)/2})$ for Hessian-free covariance estimation from an SGD
      trajectory with step exponent $a$ \citep{ni2026online}, which is at best $N^{-1/4}$ over
      the admissible range $a>1/2$: a quasi-Newton method is not in that class, because it
      already maintains curvature and forms block gradients.
\item \textbf{An exponentially small block-length bias} (\Cref{thm:lrv}). What makes
      $b\asymp\log N$ admissible is that the two-scale (flat-top / lugsail) correction
      $\Sft=\Sbm+b(\hat\Lambda_1+\hat\Lambda_1^{\!\top})$ has bias $O(\rho^b)$ rather than
      the $O(1/b)$ of plain batch means. (A flat-top lag window is one whose weights are flat
      near lag zero, so its bias picks up only the far tail of the autocovariances rather than
      shaving every lag.) The correction and this tail-sum property are due
      to the flat-top literature
      \citep{politis1995bias,politis2011higher,vats2022lugsail,singh2025utility}, not to us
      (\Cref{sec:lrv} is explicit about the boundary); what we add is that it is available
      at \emph{block-gradient} resolution on a dependent stream, and a streaming (dyadic)
      implementation that needs no advance knowledge of $N$.
\item \textbf{Exactly how badly dependence-blind intervals fail} (\Cref{prop:coverage}).
      Their asymptotic coverage is $2\Phi(z_{\alpha/2}/\sqrt{\kappa_j})-1$; over a sweep in
      dependence strength the largest gap between prediction and measurement is
      $\CovLawMaxErr$, and $\CovLawMaxErrModerate$ for $\kappa$ up to
      $\CovLawKappaModerate$, with no fitted constants (\Cref{fig:coverage}). This turns a
      previously reported qualitative failure into a quantitative one.
\item \textbf{A stability condition that blocking creates, and its fix}
      (\Cref{prop:shift}). Blocking the gradient makes the step's contraction condition
      $\gamma_t\opnorm{\Hb^{-1}\Hh_t}<2$ non-trivial, where $\gamma_t$ is the step size,
      $\Hb$ the running curvature estimate used as preconditioner, and
      $\Hh_t=b^{-1}\sum_{i\in B_t}\alpha_i\Psi_i\Psi_i^{\!\top}$ the Hessian of block $t$
      alone. On a dependent stream a block of $b$ observations can span far fewer than $b$
      distinct covariate directions, so $\Hh_t$ is near-singular and the norm is large. We give the condition, a step-size shift $t_0$ read off
      the warm-up window that restores it at $O(c_{\mathrm w}d^2)$ one-off cost, and a proof
      that the shift changes no asymptotic statement. Without it our own method diverges on a
      regime-switching design; with it that design is stable and its coverage tracks the
      infeasible oracle interval at every sample size (\Cref{tab:hard,tab:ablation}). We
      report this because it is a cost of blocking, and it is ours.
\item \textbf{Curvature schedules} (\Cref{prop:gap}). At a matched number of rank-one
      updates, deterministic gapping incurs excess curvature variance $O(\rho^m)$ whereas
      the Bernoulli thinning used as the cost knob in the independent-data method incurs
      $(\kappa_H-1)/m$. Gapping is thus a strict, if modest, improvement on the existing
      cost knob, and it makes the per-block cost deterministic.
\item \textbf{Evidence} (\Cref{sec:exp}). Four dependent designs with exact population
      targets (three with closed-form $S$), plus hourly traffic and air-quality streams
      under two protocols, one of which has an exact ground truth on real covariates.

\end{enumerate}

\paragraph{What we do not claim.}
The $O(dN)$ budget and the idea of thinning curvature updates are due to
\citet{godichonbaggioni2025adaptive}; blocked or growing mini-batches as a
dependence-breaking device are due to \citet{godichonbaggioni2023learning}; the long-run
sandwich limit under mixing or Markovian sampling is known for first-order stochastic
approximation \citep{liu2023mixing,li2023markovian,samsonov2025statistical,li2026stream};
recursive online estimation of the \emph{instantaneous} asymptotic variance under
independent data is due to \citet{godichonbaggioni2019online} and
\citet{chen2020statistical,zhu2023online}; online estimation of a \emph{long-run} covariance
under dependent sampling is known for first-order methods
\citep{roy2023online,samsonov2025statistical}; online covariance estimation for streaming
second-order methods is known under independent data
\citep{kuang2026online,wang2026inference}; the two-scale correction is known
\citep{politis1995bias,vats2022lugsail,singh2025utility}; the semiparametric efficiency
bound under Markovian sampling is due to \citet{li2023markovian}, and we attain rather
than establish it. \Cref{sec:related} states each boundary precisely.

